\subsection{张量代数中的某些运算.}

在本小节中，我们将用 \(e_x\)（\(x \in E\)）记张量代数 T(E) 到自身的线性映射 \(u \to x \otimes u\)。

引理 1. — 设 f 是 E 的对偶 E* 的一个元素。存在一个且只有一个 T(E) 到自身的线性映射 \(i_{f}\)，使得

\begin{enumerate}
\def\labelenumi{(\arabic{enumi})}
\setcounter{enumi}{3}
\tightlist
\item
\end{enumerate}

\[
i _ {f} (1) = 0\tag{5}
\]

\[
i _ {f} \circ e _ {x} + e _ {x} \circ i _ {f} = f (x). I
\]

对一切 \(x \in \mathbf{E}\)

（其中 I 表示恒等映射）。映射 \(f \to i_f\)（从 E* 到 \(\mathfrak{U}(\mathrm{T}(\mathrm{E}))\) 中）是线性的。有 \(i_f(\mathrm{T}^n) \subset \mathrm{T}^{n-1}\)，\((i_f)^2 = 0\)，且 \(i_f \circ i_g + i_g \circ i_f = 0\) 对 E* 中的 \(f\)、\(g\) 成立。映射 \(i_f\) 在子代数

T(E)（由 f 的核生成的）上为零。理想 I(Q) 在 \(i_{f}\) 下稳定；因此通过过渡到商，\(i_{f}\) 定义 C(Q) 到自身的一个线性映射（仍记作 \(i_{f}\)）。

  事实上，公式 (5) 可以写成

\[
i _ {f} (x \otimes u) = - x \otimes i _ {f} (u) + f (x). u \quad (x \in \mathrm{E}, u \in \mathrm{T(E)}).\tag{6}
\]

  由于 (4) 完全确定了 \(i_{f}\) 在 \(\mathbf{T}^{0}\) 上的值，且 (6) 确定了 \(i_{f}\) 在 \(\mathbf{T}^{n}\) 上的值（当其在 \(\mathbf{T}^{n-1}\) 上的值已知时），\(i_{f}\) 的唯一性得证。另一方面，对 \(x\in\mathbf{E}\) 与 \(u\in\mathbf{T}^{n-1}\)，(6) 的右边在 \(\mathbf{E}\times\mathbf{T}^{n-1}\) 上是双线性的；这证明了 \(i_{f}\) 的存在性（对 \(n\) 作归纳；第三章 § 1 第 2 小节），并且对 \(n\) 作归纳可知 \(i_{f}(\mathbf{T}^{n})\subset\mathbf{T}^{n-1}\)。若 \(f=ag+bh\)（\(a,b\) 属于 A，\(g,h\) 属于 E*），则显然 \(ai_{g}+bi_{h}\) 满足 (4) 与 (5)，从而等于 \(i_{f}\)。我们有

\[
(i _ {f}) ^ {2} \circ e _ {x} = - i _ {f} \circ e _ {x} \circ i _ {f} + f (x) i _ {f} = e _ {x} \circ (i _ {f}) ^ {2}
\]

  又，由于 \((i_f)^2(1) = 0\)，我们对 \(n\) 作归纳推出 \((i_f)^2\) 在 \(\mathbf{T}^n\) 上为零。把 \(f\) 换成 \(f + g\)，关系 \((i_f)^2 = 0\) 就给出 \(i_f \circ i_g + i_g \circ i_f = 0\)。一个与前面类似的归纳论证表明 \(i_f\) 在由 \(f\) 的核生成的子代数上为零。最后，(6) 蕴涵 I(Q) 中元素 \(u\)（使 \(i_f(u) \in \mathrm{I}(Q)\)）所成之集是 T(E) 的一个左理想；此外，若 \(u = (x \otimes x - Q(x).1) \otimes v (x \in E, v \in T(E))\)，则有

\[
\begin{array}{r l} i _ {f} (u) & = f (x) x \otimes v - x \otimes i _ {f} (x \otimes v) - \mathrm{Q} (x) i _ {f} (v) \\ & = f (x) x \otimes v - f (x) x \otimes v + x \otimes x \otimes i _ {f} (v) - \mathrm{Q} (x) i _ {f} (v) \\ & = (x \otimes x - \mathrm{Q} (x). 1) \otimes i _ {f} (v); \end{array}
\]

  因此 I(Q) 在 \(i_{f}\) 下稳定，最后的断言随之得出。证毕。

  设 F 是 E 上的一个双线性型；在本节余下部分，我们将用 \(i_x^{\mathrm{F}}\)（\(x \in \mathrm{E}\)）记映射 \(i_f\)，它对应于 E 上的线性型 \(f: y \to \mathrm{F}(x, y)\)。

引理 2. --- 存在 T(E) 到自身的一个线性映射 \(\lambda_{\mathrm{F}}\)，使得

\begin{enumerate}
\def\labelenumi{(\arabic{enumi})}
\setcounter{enumi}{6}
\tightlist
\item
\end{enumerate}

\[
\lambda_ {\mathrm{F}} (1) = 1\tag{8}
\]

\[
\lambda_ {\mathrm{F}} \circ e _ {x} = (e _ {x} + i _ {x} ^ {\mathrm{F}}) \circ \lambda_ {\mathrm{F}}\tag{\( (x \in E) \) .}
\]

  无论 \(f \in \mathbf{E}^*\) 如何，都有

\[
\lambda_ {\mathrm{F}} \circ i _ {f} = i _ {f} \circ \lambda_ {\mathrm{F}}.\tag{9}
\]

  事实上，公式 (8) 等价于

\[
\lambda_ {\mathbb {F}} (x \otimes u) = x \otimes \lambda_ {\mathbb {F}} (u) + i _ {x} ^ {\mathrm{F}} (\lambda_ {\mathbb {F}} (u)) \quad (x \in \mathrm{E}, u \in \mathrm{T} (\mathrm{E})).\tag{10}
\]

  由于 (7) 完全确定了 \(\lambda_{\mathbb{F}}\) 在 \(\mathrm{T}^{0}\) 上的值，且 (10) 确定了 \(\lambda_{\mathbb{F}}\) 在 \(\mathrm{T}^{n}\) 上的值（当其在 \(\mathrm{T}^{n-1}\) 上的值已知后），\(\lambda_{\mathbb{F}}\) 的唯一性得证。另一方面，对 \(x \in \mathrm{E}\) 与 \(u \in \mathrm{T}^{n-1}\)，(10) 的右边在 \(\mathrm{E} \times \mathrm{T}^{n-1}\) 上是双线性的；这证明了 \(\lambda_{\mathbb{F}}\) 的存在性（对 \(n\) 作归纳）。剩下的要证明 (9)，我们将用归纳法来做；(9) 的两边在 \(\mathrm{T}^{0}\) 上为零；设 (9) 在 \(\sum_{h=0}^{n-1} \mathrm{T}^{h}\) 上成立；于是对 \(x \in \mathrm{E}\) 与 \(u \in \mathrm{T}^{n-1}\) 我们有：

\[
\begin{array}{r l} (\lambda_ {\mathbb {F}} \circ i _ {f}) (x \otimes u) & = (- \lambda_ {\mathbb {F}} \circ e _ {x} \circ i _ {f} + f (x) \lambda_ {\mathbb {F}}) (u) \\ & = - (e _ {x} + i _ {x} ^ {\mathrm{F}}) \circ \lambda_ {\mathbb {F}} \circ i _ {f} (u) + f (x) \lambda_ {\mathbb {F}} (u) \\ & = - (e _ {x} + i _ {x} ^ {\mathrm{F}}) \circ i _ {f} \circ \lambda_ {\mathbb {F}} (u) + f (x) \lambda_ {\mathbb {F}} (u) \\ & = (i _ {f} \circ e _ {x} \circ \lambda_ {\mathbb {F}} - f (x) \lambda_ {\mathbb {F}} + i _ {f} \circ i _ {x} ^ {\mathrm{F}} \circ \lambda_ {\mathbb {F}} + f (x) \lambda_ {\mathbb {F}}) (u) \\ & = (i _ {f} \circ (e _ {x} + i _ {x} ^ {\mathrm{F}}) \circ \lambda_ {\mathbb {F}}) (u) = (i _ {f} \circ \lambda_ {\mathbb {F}}) (x \otimes u), \end{array}
\]

  从而得到我们的最后一个断言。

引理 3. --- 设 F 与 G 是 E 上两个双线性型。则我们有 \(\lambda_{\mathrm{F}} \circ \lambda_{\mathrm{G}} = \lambda_{\mathrm{F+G}}\)。对 E 上每个双线性型 F，\(\lambda_{\mathrm{F}}\) 是 T(E) 到自身上的一个双射。

  事实上，\(\lambda_{\mathbb{F}} \circ \lambda_{\mathbb{G}}\) 具有 \(\lambda_{\mathbb{F} + \mathbb{G}}: \text{有 } (\lambda_{\mathbb{F}} \circ \lambda_{\mathbb{G}})(1) = 1\) 的特征性质 (7) 与 (8)，并且

\[
\begin{array}{r l} \lambda_ {\mathrm{F}} \circ \lambda_ {\mathrm{G}} \circ e _ {x} & = \lambda_ {\mathrm{F}} \circ (e _ {x} + i _ {x} ^ {\mathrm{G}}) \circ \lambda_ {\mathrm{G}} = (e _ {x} + i _ {x} ^ {\mathrm{G}} + i _ {x} ^ {\mathrm{F}}) \circ \lambda_ {\mathrm{F}} \circ \lambda_ {\mathrm{G}} \\ & = (e _ {x} + i _ {x} ^ {\mathrm{F+G}}) \circ \lambda_ {\mathrm{F}} \circ \lambda_ {\mathrm{G}}. \end{array}
\]

  另一方面，若 F = 0，则 \(i_x^{\mathbb{F}} = 0\)（对一切 \(x \in \mathbb{E}\)），从而 \(\lambda_{\mathbb{F}} = \mathbb{I}\)，这蕴涵对每个 F 有 \(\lambda_{\mathbb{P}} \circ \lambda_{-\mathbb{P}} = \lambda_{-\mathbb{F}} \circ \lambda_{\mathbb{F}} = \mathbb{I}\)。

